% Research contribution: constructive gamma identities.
% This is a section fragment, to be incorporated into the main article.
% All limiting language below abbreviates the displayed finite error estimates.

\section{Finite trigonometry and Euler's reflection formula}

\subsection{The elementary trigonometric foundation}

Here \(\exp\), \(\sin\), and \(\cos\) are the functions given by their
factorial power series, with the explicitly bounded tails described earlier.
Put
\[
 A(t)=\int_0^t\frac{du}{1+u^2},\qquad \pi=4A(1).
\]
The alternating expansion
\[
 A(1)=\sum_{j=0}^{L-1}\frac{(-1)^j}{2j+1}+r_L,
 \qquad |r_L|\le\frac1{2L+1}
\]
is obtained by integrating the finite geometric identity for
\((1+u^2)^{-1}\). Thus \(\pi\) has a primitive recursive rational name.
Finite rational comparisons of the first eight terms give \(3<\pi<4\).
Multiplication of finite exponential series, followed by their factorial
tail bounds, proves the addition laws and
\(\exp(it)=\cos t+i\sin t\).

There is no hidden appeal here to the qualitative theory of trigonometric
functions. For \(0\le t\le1\), differentiate
\[
 e^{-2iA(t)}\frac{1+it}{1-it}.
\]
The two terms cancel exactly. The finite fundamental theorem of calculus
therefore gives
\[
 e^{2iA(t)}=\frac{1+it}{1-it},\qquad
 \sin(2A(t))=\frac{2t}{1+t^2},\qquad
 \cos(2A(t))=\frac{1-t^2}{1+t^2}.                 \tag{T1}
\]
In particular, \(e^{i\pi/2}=i\), \(e^{i\pi}=-1\), and
\(e^{2\pi i}=1\). Since \((2A)'(t)=2/(1+t^2)\ge1\), bounded
bisection constructs its inverse on \([0,\pi/2]\), with an explicit
linear modulus. Formula (T1) proves positivity of sine there. Its second
derivative is \(-\sin\), so the finite second-derivative comparison with
the chord joining the endpoints gives
\[
 \sin v\ge\frac{2v}{\pi}\quad(0\le v\le\pi/2).
                                                        \tag{T2}
\]
For completeness, the chord comparison follows by applying the finite
mean-value inequality to the derivative: a differentiable function with
nonpositive second derivative has decreasing secant slopes. Both this
assertion and (T1) can instead be checked on rational meshes; their
derivative bounds give an explicit mesh error. No unrestricted intermediate
value theorem is used: the only inverse just mentioned has slope at least
one and is obtained by a prescribed bisection program. This program does
not compare exact real values. If the current bracket has width \(w\),
approximate the midpoint value minus the target to error \(w/16\).
When the rational approximation exceeds \(w/8\) in absolute value,
its sign selects a half bracket. Otherwise the slope bound puts the
inverse in the interval of radius \(w/4\) about the midpoint. Intersect
this interval with the old bracket. In every case the width halves, and
every decision is a rational comparison.

\subsection{A finite sine factorization}

Define the entire function
\[
 \operatorname{sinc}(w)=\sum_{j=0}^{\infty}
                \frac{(-1)^jw^{2j}}{(2j+1)!}.
\]
Thus \(\sin w=w\operatorname{sinc}(w)\), including at \(w=0\).
For an odd integer \(m=2J+1\), repeated use of the addition law writes
\(\sin(mu)/(m\sin u)\) as a polynomial of degree \(J\) in
\(\sin^2u\), with constant term one. Its explicitly given distinct roots
are \(\sin^2(k\pi/m)\), \(1\le k\le J\). Consequently the following is
a finite polynomial identity, not an identity obtained by analytic
continuation:
\[
 \operatorname{sinc}(mu)=\operatorname{sinc}(u)
 \prod_{k=1}^{J}
 \left(1-\frac{\sin^2u}{\sin^2(k\pi/m)}\right).           \tag{T3}
\]
To justify the polynomial argument constructively, use the standard
multiple-angle recurrence (equivalently, the recurrence for the
polynomials \(U_j\)). For odd \(m\), \(U_{m-1}(\cos u)\) is a polynomial
in \(1-\sin^2u\). The listed roots are obtained by substituting the
angles \(k\pi/m\), and their distinctness follows from positivity of
cosine on this finite interval. More quantitatively,
\(\cos v\ge1/m\) for
\(0\le v\le(m-1)\pi/(2m)\), by (T2) applied to the complementary
angle. Consecutive sine values therefore differ by at least
\(\pi/m^2\). Repeated monic polynomial division by the known linear
factors proves the factorization; no fundamental theorem of algebra is
needed. Multiplication by \(\operatorname{sinc}(u)\) removes all
division by a possibly zero \(\sin u\).

\subsection{An explicit sine-product error bound}

Let \(R\ge1\) and \(K\ge1\) be integers, let \(|z|\le R\), and put
\[
 m=8K^2R^2+1,\quad J=(m-1)/2,\quad
 S_K(z)=\prod_{k=1}^K\left(1-\frac{z^2}{k^2}\right).
\]
Then
\[
 \left|\operatorname{sinc}(\pi z)-S_K(z)\right|
 \le\frac{2^{74R^2}(R+1)^2}{K}.                         \tag{T4}
\]
Here is a complete estimate. Write
\(u=\pi z/m\), \(v_k=\pi k/m\), and
\(a_k=\sin^2u/\sin^2v_k\). Since
\(\sin v_k\ge2k/m\) and \(|\sin u|\le2|u|\),
\[
 |a_k|\le\frac{16R^2}{k^2},\qquad
 \sum_{k=1}^J|a_k|\le32R^2,
 \qquad \sum_{k=1}^K|z^2/k^2|\le2R^2.                  \tag{T5}
\]
For \(|w|\le1\), the factorial series gives
\(|\operatorname{sinc}(w)-1|\le|w|^2/2\) and
\(|\operatorname{sinc}(w)|\le2\). For \(k\le K\), (T2) also gives
\(\operatorname{sinc}(v_k)>1/2\). Thus
\[
 \left|a_k-\frac{z^2}{k^2}\right|
 \le\frac{96R^2}{m^2}\left(\frac{R^2}{k^2}+1\right).
\]
Indeed, \(|\operatorname{sinc}(u)^2-1|\le3|u|^2/2\), while
\(0<\operatorname{sinc}(v_k)\le1\) implies
\(|1-\operatorname{sinc}(v_k)^2|\le v_k^2\). Dividing their sum by
\(\operatorname{sinc}(v_k)^2>1/4\) bounds the ratio difference by
\(6(|u|^2+v_k^2)\); now use \(\pi^2<16\).
The tail satisfies \(\sum_{k>K}|a_k|\le16R^2/K\). The finite
telescoping identity for a difference of products, together with
\(1+t\le e^t\), therefore yields
\[
 \left|\prod_{k=1}^J(1-a_k)-S_K(z)\right|
 \le e^{34R^2}\left\{\frac{16R^2}{K}
       +\frac{96R^2(2R^2+K)}{m^2}\right\}.              \tag{T6}
\]
By (T3), one must add at most
\(8R^2e^{32R^2}/m^2\) to account for the prefactor
\(\operatorname{sinc}(u)\). Substituting the displayed value of \(m\)
bounds the braces, including this extra term, by \(32R^2/K\).
Finally \(e^{34R^2}\le2^{68R^2}\), which proves the deliberately
generous bound (T4).

This proves, in the precise effective sense needed here,
\[
 \frac{\sin\pi z}{\pi z}
       =\prod_{k=1}^{\infty}\left(1-\frac{z^2}{k^2}\right),             \tag{T7}
\]
where the left side means \(\operatorname{sinc}(\pi z)\) at every
point. For an error \(2^{-p}\), it is enough to take
\(K\ge2^{p+74R^2}(R+1)^2\). Every product and every trigonometric
value occurring in this argument is evaluated to a prescribed rational
precision; its finite error budget is obtained from (T5)--(T6).

\subsection{Reflection and constructive division}

Use the Euler approximants
\[
 \Gamma_N(z)=\frac{N!\,N^z}{z(z+1)\cdots(z+N)}.
\]
At points with the displayed denominator factors apart from zero, exact
finite cancellation gives
\[
 \Gamma_N(z)\Gamma_N(-z)
       =-\frac1{z^2S_N(z)}.                              \tag{R1}
\]
The established convergence of these approximants, the recurrence
\(\Gamma(1-z)=-z\Gamma(-z)\), and (T4) therefore prove
\[
 \boxed{\ \Gamma(z)\Gamma(1-z)=\frac\pi{\sin\pi z}\ }.              \tag{R2}
\]
The domain assertion is constructive: supply a radius \(|z|\le R\)
and a positive rational \(\delta\le1\) separating \(z\) from the
finitely many integers in \([-2R-1,2R+1]\). Integers farther away have
an automatic separation. This is stronger information than the bare
negation ``\(z\) is not an integer,'' and it is exactly what division
requires.

One may extract an explicit lower bound for the sine denominator.
Choose an integer \(L\ge2R\). The first \(L\) factors in (T7) have
product of moduli at least \(\delta^{2L}/(L!)^2\). For the remaining
factors, \(|z^2/k^2|\le1/4\), and the real inequality
\(\log(1-t)\ge-2t\) gives a product lower bound
\(e^{-2R^2/L}\ge2^{-4R^2}\). Hence
\[
 |\sin\pi z|\ge
       \frac{\delta^{2L+1}2^{-4R^2}}{(L!)^2}.             \tag{R3}
\]
This coarse rational bound makes every passage to a reciprocal in (R2)
primitive recursive. Alternatively, the cleared identity
\(\sin(\pi z)\Gamma(z)\Gamma(1-z)=\pi\) can be proved first.

At \(z=1/2\), positivity of the defining gamma integral selects the
positive root and gives
\[
 \boxed{\Gamma(1/2)=\sqrt\pi}.                           \tag{R4}
\]
The square root here is the positive one, computed by bounded rational
bisection. No version of Stirling's formula has been used.

\section{Gauss multiplication without an asymptotic normalization}

\subsection{A second finite sine identity}

For every integer \(m\ge1\),
\[
 \prod_{r=1}^{m-1}\sin\frac{\pi r}{m}
                 =\frac{m}{2^{m-1}}.                    \tag{M1}
\]
For \(m>1\), put \(\omega=e^{2\pi i/m}\). The explicitly known
numbers \(1,\omega,\ldots,\omega^{m-1}\) are pairwise apart: (T2)
gives \(|\omega^r-\omega^s|\ge4/m\) when \(r\ne s\) in this range.
Polynomial division by these known roots proves
\(X^m-1=\prod_{r=0}^{m-1}(X-\omega^r)\). Divide out \(X-1\) and
substitute \(X=1\). Taking squared moduli gives
\[
 m^2=\prod_{r=1}^{m-1}4\sin^2(\pi r/m).
\]
Every sine factor is positive, by (T1), reflection about \(\pi/2\),
and (T2); positive square-root cancellation proves (M1). For \(m=1\)
the empty product is one.

\subsection{Exact grouping of Euler approximants}

Fix \(m\ge1\), set \(a=1/m\), and for \(n\ge1\) set
\(L=m(n+1)-1\). Denominator grouping gives the exact identity
\[
 U_n(z):=\frac{\prod_{r=0}^{m-1}\Gamma_n(z+r/m)}{\Gamma_L(mz)}
   =\frac{(n!)^m m^{m(n+1)}}{L!}\,
       n^{(m-1)/2}\left(\frac nL\right)^{mz}.
\]
All factors independent of \(z\) disappear when this expression is
divided by its value at \(a\):
\[
 \frac{U_n(z)}{U_n(a)}
                 =\left(\frac nL\right)^{m(z-a)}.        \tag{M2}
\]
Only powers of positive real numbers occur, so no branch choice is
implicit. Moreover,
\[
 \left|\log\frac{mn}{L}\right|
   =\log\left(1+\frac{m-1}{mn}\right)\le\frac1n.
\]
For \(|z|\le R\) and \(n\ge m(R+1)\), this proves the explicit
estimate
\[
 \left|\left(\frac nL\right)^{m(z-a)}
           -m^{-m(z-a)}\right|
 \le\frac{2m(R+1)m^{m(R+1)}}n.                           \tag{M3}
\]
Consequently convergence of the gamma approximants proves
\[
 \frac{\prod_{r=0}^{m-1}\Gamma(z+r/m)}{\Gamma(mz)}
   = m^{1-mz} P_m,
 \qquad P_m=\prod_{r=1}^{m-1}\Gamma(r/m).                \tag{M4}
\]
One supplies a positive rational separation of \(mz\) from the
nonpositive integers in the relevant finite disk. This also supplies
separations for all \(z+r/m\): multiplication by \(m\) sends their
possible poles to nonpositive integers. The reciprocal-product bounds
for gamma handle the divisions in (M2)--(M4).

The constant is found without any asymptotic theorem. Reflection and
(M1) give
\[
 P_m^2=\prod_{r=1}^{m-1}
          \Gamma(r/m)\Gamma(1-r/m)
       =\frac{\pi^{m-1}}{\prod_{r=1}^{m-1}\sin(\pi r/m)}
       =\frac{(2\pi)^{m-1}}m.
\]
Since \(P_m>0\), substitution in (M4) proves
\[
 \boxed{\ \prod_{r=0}^{m-1}\Gamma(z+r/m)
       =(2\pi)^{(m-1)/2}m^{1/2-mz}\Gamma(mz)\ }.        \tag{M5}
\]
In particular,
\[
 \Gamma(z)\Gamma(z+1/2)
       =2^{1-2z}\sqrt\pi\,\Gamma(2z),                   \tag{M6}
\]
and
\[
 \Gamma(z)\Gamma(z+1/3)\Gamma(z+2/3)
       =2\pi\,3^{1/2-3z}\Gamma(3z).                    \tag{M7}
\]
Thus duplication is available to determine the constant in Stirling's
formula without circular reasoning.

\section{Beta integrals from one-dimensional finite estimates}

\subsection{Definition, tails, and recurrences}

For certified \(\Re a\ge\alpha>0\), \(\Re b\ge\beta>0\), define
\[
 B(a,b)=\int_0^1t^{a-1}(1-t)^{b-1}\,dt.                 \tag{B1}
\]
All powers use the real logarithm of the positive variable. On
\([\epsilon,1-\epsilon]\), compute a finite rational Riemann sum,
with the explicit derivative bound obtained by differentiating the
integrand. For \(0<\epsilon\le1/2\), the omitted endpoints have
absolute integral at most
\[
 \frac{2\epsilon^\alpha}{\alpha}
       +\frac{2\epsilon^\beta}{\beta}.                  \tag{B2}
\]
This defines a primitive recursive name and a Cauchy modulus; small
rational \(\epsilon=2^{-h}\) is chosen with the displayed bound below
the requested error. Symmetry is the finite substitution \(t\mapsto1-t\).
The identity \(1=t+(1-t)\), and integration of the derivative of
\(t^a(1-t)^b\), give
\[
 B(a,b)=B(a+1,b)+B(a,b+1),\qquad
 aB(a,b+1)=bB(a+1,b).
\]
Indeed, on the truncated interval these are finite integration-by-parts
identities up to the mesh error; the boundary terms have modulus at most
\(\epsilon^\alpha+\epsilon^\beta\). Hence
\[
 B(a+1,b)=\frac a{a+b}B(a,b),\qquad
 B(a,b+1)=\frac b{a+b}B(a,b).                            \tag{B3}
\]
The denominator has real part at least \(\alpha+\beta\), so division
has a supplied margin. Iterating gives
\[
 B(a,b+N)=\frac{(b)_N}{(a+b)_N}B(a,b),                  \tag{B4}
\]
where \((w)_N=\prod_{j=0}^{N-1}(w+j)\). In particular,
\[
 B(a,N+1)=\frac{N!}{a(a+1)\cdots(a+N)}.                \tag{B5}
\]

\subsection{An effective beta-to-gamma limit}

For fixed bounded \(a,b\) as above,
\[
 N^aB(a,b+N)\longrightarrow\Gamma(a).                   \tag{B6}
\]
Here is an estimate sufficient to implement the assertion. Choose
integers \(d\ge\max(1,\Re a)\) and \(C\ge|b-1|\), and write
\(H=1/\alpha+d!\). For \(T\ge1\),
\(N\ge\max(2T,2,T^2+2CT)\), one has
\[
 |N^aB(a,b+N)-\Gamma(a)|
 \le \frac{2H(T^2+2CT)}N
         +2\,4^{d+1}d!\,e^{-T/4}.                     \tag{B7}
\]
To prove it, put \(t=s/N\) in (B1), obtaining an integral on
\([0,N]\). On \([0,T]\), set
\[
 D(s)=(N+b-1)\log(1-s/N)+s.
\]
The elementary inequalities
\(-\log(1-u)-u\le u^2\) and \(|\log(1-u)|\le2u\) for
\(0\le u\le1/2\) yield
\(|D(s)|\le(T^2+2CT)/N\le1\). Therefore
\(|e^{D(s)}-1|\le2|D(s)|\). Since
\(\int_0^\infty s^{\Re a-1}e^{-s}\,ds\le H\), this proves the first
term in (B7). On the remaining interval, \(N+\Re b-1\ge N/2\), so
\[
 |(1-s/N)^{N+b-1}|\le e^{-s/2}.
\]
For \(s\ge T\ge1\), use
\(s^{\Re a-1}\le s^d\) and
\(s^de^{-s/2}\le4^dd!e^{-s/4}\). Integration bounds each of the two
tails by \(4^{d+1}d!e^{-T/4}\), proving (B7).

For example, for requested error \(2^{-p}\), one may choose
\[
 T=4(p+d^2+2d+5),\quad
 N\ge\max\{2T,2,T^2+2CT,
       2^{p+3}(T^2+2CT)\lceil H\rceil\}.
\]
The elementary estimates \(d!\le2^{d^2}\) and \(e\ge2\) verify the
error directly. There is no invocation of dominated convergence.

For \(N\ge2\), the exact Euler-approximant relation is
\[
 N^a\frac{(b)_N}{(a+b)_N}
   =\left(\frac N{N-1}\right)^a
                 \frac{\Gamma_{N-1}(a+b)}{\Gamma_{N-1}(b)}.
\]
Use this in (B4), then apply (B7) and the effective convergence of
\(\Gamma_N\). One obtains first the cleared identity, and then, by the
known nonvanishing margin for gamma,
\[
 \boxed{\ B(a,b)=\frac{\Gamma(a)\Gamma(b)}{\Gamma(a+b)}\ }.          \tag{B8}
\]
This proof uses only one-dimensional substitutions and finite products;
a general change-of-variables theorem for improper double integrals is
unnecessary.

\subsection{Consequences and integral evaluations}

The substitution \(t=u/(1+u)\) in truncated intervals proves
\[
 \int_0^\infty\frac{u^{a-1}}{(1+u)^{a+b}}\,du=B(a,b). \tag{B9}
\]
The omitted tails have bounds \(\epsilon^\alpha/\alpha\) at zero and
\(R^{-\beta}/\beta\) at infinity. The substitution
\(t=\sin^2\theta\), again first on a closed subinterval, gives
\[
 \int_0^{\pi/2}\sin^{2a-1}\theta\cos^{2b-1}\theta\,d\theta
                     =\tfrac12 B(a,b).                \tag{B10}
\]
The beta endpoint estimates (B2) control the removed angular intervals.
For natural numbers \(j,k\),
\[
 B(j+1,k+1)=\frac{j!k!}{(j+k+1)!}.                      \tag{B11}
\]
Combining (B8)--(B9) with reflection gives, for certified
\(0<\Re a<1\),
\[
 \int_0^\infty\frac{t^{a-1}}{1+t}\,dt
       =\frac\pi{\sin\pi a}.                          \tag{B12}
\]
For positive rational \(q\) and \(0<\Re a<q\), replacing \(t\) by
\(u^{1/q}\) gives
\[
 \int_0^\infty\frac{t^{a-1}}{1+t^q}\,dt
       =\frac\pi{q\sin(\pi a/q)}.                     \tag{B13}
\]
Similarly, for positive rationals \(c,q\),
\[
 \int_0^\infty t^{a-1}e^{-ct^q}\,dt
       =\frac{\Gamma(a/q)}{q\,c^{a/q}},\qquad \Re a>0. \tag{B14}
\]
All substitutions are between explicitly truncated intervals; their
endpoint errors follow from (B2), (B7), or the gamma tail estimates.

\section{Special values and exact rational-argument relations}

Recurrence and (R4) give, for \(n\ge0\),
\[
 \Gamma(n+1/2)=\frac{(2n)!}{4^n n!}\sqrt\pi,\qquad
 \Gamma(1/2-n)=\frac{(-4)^n n!}{(2n)!}\sqrt\pi.          \tag{V1}
\]
The finite verification is induction: each additional factor is
\(n+1/2\) in the first formula and \(1/2-n\) in the second.
Thus, for example, \(\Gamma(-1/2)=-2\sqrt\pi\) and
\(\Gamma(-3/2)=4\sqrt\pi/3\).

Reflection supplies
\[
 \Gamma(1/3)\Gamma(2/3)=\frac{2\pi}{\sqrt3},\quad
 \Gamma(1/4)\Gamma(3/4)=\pi\sqrt2,\quad
 \Gamma(1/6)\Gamma(5/6)=2\pi.                           \tag{V2}
\]
Let \(A=\Gamma(1/3)\), \(D=\Gamma(1/4)\), both positive. Duplication
at \(1/6\), followed by (V2), gives
\[
 \Gamma(1/6)=\frac{2^{-1/3}\sqrt3}{\sqrt\pi}A^2,
 \qquad
 \Gamma(5/6)=\frac{2^{4/3}\pi^{3/2}}{\sqrt3\,A^2}.     \tag{V3}
\]
Duplication and triplication also give the useful twelfth-argument
relations
\[
 \Gamma(1/12)\Gamma(7/12)=\sqrt6\,A^2,
 \qquad
 \Gamma(1/12)\Gamma(5/12)=\sqrt2\,3^{1/4}D^2.           \tag{V4}
\]
Combining these with reflection proves the explicit positive-square
identities
\[
 \Gamma(1/12)^2
   =\frac{3^{3/4}(\sqrt6+\sqrt2)}{2\pi}A^2D^2,
 \qquad
 \Gamma(5/12)^2
   =\frac{\pi(\sqrt6-\sqrt2)}{3^{1/4}}\frac{D^2}{A^2}. \tag{V5}
\]
All roots in these expressions are positive roots. They are computed by
bounded bisection, and their algebraic identities are verified by
polynomial multiplication together with positive lower bounds.

For fifths, reflection yields
\[
 \Gamma(1/5)\Gamma(4/5)
       =\frac{4\pi}{\sqrt{10-2\sqrt5}},\qquad
 \Gamma(2/5)\Gamma(3/5)
       =\frac{4\pi}{\sqrt{10+2\sqrt5}}.                 \tag{V6}
\]
More generally, the finite normalization computation already proved
\[
 \prod_{r=1}^{m-1}\Gamma(r/m)
          =\frac{(2\pi)^{(m-1)/2}}{\sqrt m}.             \tag{V7}
\]
These are relations among values, not claims that all rational-argument
gamma values reduce to elementary constants. A finite certificate for a
proposed relation can consist of applications of recurrence, reflection,
and multiplication, followed by exact rational and algebraic arithmetic.

For the rising factorial, the recurrence also proves
\[
 \Gamma(z+n)=(z)_n\Gamma(z),                            \tag{V8}
\]
at arguments with the requisite pole margins. For integer \(n\ge0\),
the generalized binomial polynomial satisfies
\[
 \binom zn=\frac{\Gamma(z+1)}{n!\,\Gamma(z-n+1)}        \tag{V9}
\]
whenever the displayed gamma values are defined. The left side is the
finite polynomial \(z(z-1)\cdots(z-n+1)/n!\). At an apparent cancelled
pole, use this polynomial or a reciprocal-gamma version, rather than
performing undefined division.

\section{Modulus formulas on vertical lines}

Complex conjugation commutes with every Euler approximant and hence
with gamma. Apply reflection to \(z=iy\), with real \(y\) apart from
zero, and use \(\sin(i\pi y)=i\sinh(\pi y)\). Recurrence gives
\[
 |\Gamma(iy)|^2=\frac\pi{y\sinh(\pi y)},\qquad
 |\Gamma(1+iy)|^2=\frac{\pi y}{\sinh(\pi y)}.           \tag{VM1}
\]
For all real \(y\), including zero, a uniformly constructive version is
\[
 |\Gamma(1+iy)|^2
  =\left(\sum_{j=0}^{\infty}
                \frac{(\pi y)^{2j}}{(2j+1)!}\right)^{-1}.             \tag{VM2}
\]
The denominator is at least one. To verify (VM2) without cancelling
\(y\), use the exact finite identity
\[
 \Gamma_N(1+iy)\Gamma_N(1-iy)
   =\frac{N^2}{(N+1)^2}
     \left(\prod_{k=1}^{N+1}(1+y^2/k^2)\right)^{-1}.
\]
The entire sine product at \(z=iy\) identifies the limit of this
positive product with the denominator in (VM2). Every reciprocal has
lower bound one before inversion, so no decision whether \(y=0\) is
required.
At \(z=1/2+iy\), reflection instead gives
\[
 |\Gamma(1/2+iy)|^2=\frac\pi{\cosh(\pi y)}              \tag{VM3}
\]
for every real \(y\); this denominator too is bounded below by one.
Repeated recurrence yields
\[
 |\Gamma(n+iy)|^2
  =\prod_{k=1}^{n-1}(k^2+y^2)\,|\Gamma(1+iy)|^2
       \quad(n\ge1),                                  \tag{VM4}
\]
and
\[
 |\Gamma(n+1/2+iy)|^2
  =\prod_{k=0}^{n-1}((k+1/2)^2+y^2)
                   \frac\pi{\cosh(\pi y)}\quad(n\ge0). \tag{VM5}
\]

\section{Differentiated functional identities}

The derivative estimates for the reciprocal product justify
differentiating the finite identities above with their controlled
remainders. With \(\psi=\Gamma'/\Gamma\), one obtains
\[
 \psi(1-z)-\psi(z)=\pi\cot(\pi z),\qquad
 \psi'(1-z)+\psi'(z)=\pi^2\csc^2(\pi z),                \tag{D1}
\]
and
\[
 \sum_{r=0}^{m-1}\psi(z+r/m)=m\psi(mz)-m\log m.         \tag{D2}
\]
For \(j\ge1\), further differentiation gives
\[
 \sum_{r=0}^{m-1}\psi^{(j)}(z+r/m)
                         =m^{j+1}\psi^{(j)}(mz).        \tag{D3}
\]
Combining these with \(\psi(1)=-\gamma\) yields
\[
 \begin{aligned}
 \psi(1/2)&=-\gamma-2\log2,\\
 \psi(1/3)&=-\gamma-\frac{\pi}{2\sqrt3}-\frac32\log3,\quad
 \psi(2/3)=-\gamma+\frac{\pi}{2\sqrt3}-\frac32\log3,\\
 \psi(1/4)&=-\gamma-\frac\pi2-3\log2,\quad
 \psi(3/4)=-\gamma+\frac\pi2-3\log2,\\
 \psi(1/6)&=-\gamma-\frac{\pi\sqrt3}{2}-\frac32\log3-2\log2,\\
 \psi(5/6)&=-\gamma+\frac{\pi\sqrt3}{2}-\frac32\log3-2\log2.
 \end{aligned}                                        \tag{D4}
\]
For example, reflection at \(1/3\) determines the difference of the
two third-argument values, while (D2) with \(m=3,z=1/3\) determines
their sum. The quarter- and sixth-argument formulas follow in exactly
the same finite way from duplication and reflection. Thus (D4) involves
only solving two linear equations with nonzero integer coefficients.
At \(z=1/2\), (D1) gives \(\psi'(1/2)=\pi^2/2\).

\subsection{How the identities are read in iPRA}

For each fixed collection of primitive recursive input names and
certificates, the identities in this section mean equality of their
rational approximants at every prescribed accuracy. Bounds (T4), (R3),
(M3), (B2), and (B7), together with the gamma-product bounds, determine
all truncation lengths and arithmetic precisions by primitive recursive
formulas. Product regrouping, polynomial division by the explicitly
given roots, integration by parts on a finite mesh, and all endpoint
comparisons have finite coded invariants. Quantifier-free induction
proves those invariants. The language does not quantify over an
uncertified universal evaluator for arbitrary primitive recursive
programs. The claims are uniform finite arithmetic lemmas and their
instances for the stated names, rather than an implicit appeal to a
higher-order theory of all continuous functions.

% Primary formula references, checked 2026-11-15:
% NIST DLMF 5.5, https://dlmf.nist.gov/5.5 : recurrence, reflection,
% duplication, multiplication. Retrieval reference turn3view0.
% NIST DLMF 5.12, https://dlmf.nist.gov/5.12 : beta integrals.
% Retrieval reference turn3view1.
% The finite proofs and intentionally coarse moduli above are supplied
% here independently; citations identify the standard identities only.
